% Презентація етапу 3 «Планування».
% Збирання: make pdf (слайди), make notes (слайди + нотатки на другому екрані),
% make speech (лише текст виступу, для друку).
%
% Хто говорить — у нотатках, у квадратних дужках. Розподіл за Team-Tasks.md:
%   [Лід · Яна]   — вступ, порядок на два тижні, фінал
%   [П1]          — Project Overview (продукт, цілі, обсяг, обмеження, ризики)
%   [П2 · Дарина] — Roadmap і діаграма Ґанта (автор Roadmap.md)
%   [П4 · Ярина]  — RACI і матриця Ейзенхауера (автор RACI.md)
%   [П3]          — беклог, ресурс MVP, критичний шлях
%   [П5]          — обсяг робіт на модуль 2 (автор Semester_Scope.md)
% TODO: замінити П1, П3, П5 на імена.
%
% Усі цифри — з артефактів етапу 3 (Project-Overview.md, Roadmap.md, gantt.pdf,
% RACI.md, Backlog.pdf, Eisenhower.pdf, Semester_Scope.md) і Product Map.

\ifdefined\notesmode\PassOptionsToClass{handout}{beamer}\fi
\documentclass[aspectratio=169,11pt]{beamer}
\usetheme{lapcoach}

\ifdefined\notesmode
  \ifx\notesmode\notesonly
    \setbeameroption{show only notes}
  \else
    \usepackage{pgfpages}
    \setbeameroption{show notes on second screen=right}
  \fi
\fi

\usepackage[ukrainian]{babel}
\usepackage{tabularx}
\usepackage{multirow}
% Допоміжні макроси для слайдів (винесено з кадрів, бо # усередині frame ламається)
% Блок фіксованого розміру з текстом угорі
\newcommand{\topbox}[3]{\parbox[t][#2][t]{#1}{\raggedright #3}}

% з slides/02-product.tex
\newcommand{\phasecol}[3]{%
  \begin{tikzpicture}
    \node[anchor=north west,inner sep=5pt,fill=soft,font=\tiny,text=body] (p) at (0,0)
      {\topbox{3.05cm}{3.0cm}{{\scriptsize\semibold\color{ink}#2\par}\vskip3pt #3\par}};
    \fill[#1] (p.north west) rectangle ($(p.north east)+(0,0.07)$);
  \end{tikzpicture}}

% з slides/02-product.tex
\newcommand{\limcard}[2]{%
  \begin{tikzpicture}
    \node[anchor=north west,inner sep=5pt,fill=soft] (q) at (0,0)
      {\topbox{3.0cm}{1.45cm}{{\footnotesize\semibold\color{ink}#1\par}\vskip3pt{\tiny\color{body}#2\par}}};
    \fill[kerb] (q.north west) rectangle ($(q.south west)+(0.05,0)$);
  \end{tikzpicture}}

% з slides/03-work.tex
\newcommand{\gphase}[4]{% початок, кінець, назва, номер
  \fill[ink] (#1-1+0.03,-0.36) rectangle (#2-0.03,-0.92);
  \node[anchor=west,font=\tiny\semibold,text=white,text width=(#2-#1+1)*1.183cm-0.25cm,align=left,inner sep=0pt]
    at (#1-1+0.1,-0.64) {Ф#4\enspace#3};}

% з slides/03-work.tex
\newcommand{\gwork}[6]{% фаза-початок, фаза-кінець, початок, кінець, рядок, назва
  \pgfmathsetmacro{\yy}{-1.06-(#5-1)*0.76}
  \fill[kerb] (#3-1+0.03,\yy) rectangle (#4-0.03,\yy-0.12);
  \node[anchor=north west,font=\tiny,text=body,text width=(#2-#1+1)*1.183cm-0.2cm,align=left,inner sep=0pt]
    at (#1-1+0.06,\yy-0.17) {#6};}

% з slides/03-work.tex
\newcommand{\racirow}[4]{% рядок, ім'я, роль, клітинки
  \pgfmathsetmacro{\yy}{-(#1-1)*0.62}
  \node[anchor=west,font=\scriptsize\semibold,text=ink,inner sep=0pt] at (-2.75,\yy+0.09) {#2};
  \node[anchor=west,font=\tiny,text=mute,inner sep=0pt] at (-2.75,\yy-0.16) {#3};
  \foreach \v [count=\c] in {#4}{%
    \def\cx{(\c-1)*0.42+0.21}
    \ifnum\strcmp{\v}{AR}=0
      \fill[kerb] ({\cx-0.19},\yy-0.25) rectangle ({\cx+0.19},\yy+0.25);
      \node[font=\fontsize{4.6}{5}\selectfont\semibold,text=white] at ({\cx},\yy) {A/R};
    \else\ifnum\strcmp{\v}{R}=0
      \fill[ink] ({\cx-0.19},\yy-0.25) rectangle ({\cx+0.19},\yy+0.25);
      \node[font=\tiny\semibold,text=white] at ({\cx},\yy) {R};
    \else\ifnum\strcmp{\v}{A}=0
      \draw[kerb,line width=0.8pt] ({\cx-0.17},\yy-0.23) rectangle ({\cx+0.17},\yy+0.23);
      \node[font=\tiny\semibold,text=kerb] at ({\cx},\yy) {A};
    \else\ifnum\strcmp{\v}{C}=0
      \fill[line] ({\cx-0.19},\yy-0.25) rectangle ({\cx+0.19},\yy+0.25);
      \node[font=\tiny,text=body] at ({\cx},\yy) {C};
    \else\ifnum\strcmp{\v}{I}=0
      \node[font=\tiny,text=later] at ({\cx},\yy) {I};
    \fi\fi\fi\fi\fi
  }}

% з slides/03-work.tex
\newcommand{\loadrow}[3]{%
  \pgfmathsetmacro{\yy}{-(#1-1)*0.62}
  \node[anchor=west,font=\scriptsize\semibold,text=ink,inner sep=0pt] at (-2.3,\yy) {#2};
  \foreach \v [count=\c] in {#3}{%
    \ifcase\v \def\fc{soft}\def\tc{later}\or \def\fc{kerb!30}\def\tc{ink}\or \def\fc{kerb}\def\tc{white}\fi
    \fill[\fc] ({(\c-1)*0.62+0.03},\yy-0.28) rectangle ({\c*0.62-0.03},\yy+0.28);
    \node[font=\scriptsize\semibold,text=\tc] at ({(\c-0.5)*0.62},\yy) {\v};
  }}

% з slides/03-work.tex
\newcommand{\ok}{\tikz[baseline=-0.5ex]{\fill[go] (0,0) circle (0.11); \draw[white,line width=0.9pt] (-0.06,0) -- (-0.01,-0.05) -- (0.06,0.05);}\enspace}

% з slides/04-backlog.tex
\newcommand{\brow}[7]{% рядок, фаза, період, Must, Should, Won't, години
  \pgfmathsetmacro{\yy}{-(#1-1)*0.78}
  \node[anchor=west,font=\footnotesize\semibold,text=ink,inner sep=0pt] at (0,\yy+0.1) {#2};
  \node[anchor=west,font=\tiny,text=mute,inner sep=0pt] at (0,\yy-0.17) {#3};
  \pgfmathsetmacro{\a}{2.9}
  \pgfmathsetmacro{\b}{\a+#4*0.36}
  \pgfmathsetmacro{\c}{\b+#5*0.36}
  \pgfmathsetmacro{\d}{\c+#6*0.36}
  \ifnum#4>0 \fill[kerb] (\a,\yy-0.22) rectangle (\b-0.03,\yy+0.22);
    \node[font=\tiny\semibold,text=white] at ({(\a+\b)/2},\yy) {#4};\fi
  \ifnum#5>0 \fill[kerb!35] (\b,\yy-0.22) rectangle (\c-0.03,\yy+0.22);
    \node[font=\tiny\semibold,text=ink] at ({(\b+\c)/2},\yy) {#5};\fi
  \ifnum#6>0 \fill[line] (\c,\yy-0.22) rectangle (\d-0.03,\yy+0.22);
    \node[font=\tiny\semibold,text=mute] at ({(\c+\d)/2},\yy) {#6};\fi
  \node[anchor=east,font=\footnotesize\heavy,text=ink] at (14.2,\yy) {#7};}

% з slides/04-backlog.tex
\newcommand{\hb}{0.085}% см на годину

% з slides/04-backlog.tex
\newcommand{\qcell}[7]{% x, y, колір, назва, кількість, підпис, приклади
  \node[anchor=north west,inner sep=0.15cm,fill=#3] at (#1,#2) {\topbox{4.1cm}{2.1cm}{%
    {\scriptsize\semibold\color{\qtext}#4}\hfill{\fontsize{16}{16}\selectfont\heavy\color{\qtext}#5}\par
    {\tiny\color{\qsub}#6\par}\vskip0.08cm
    {\tiny\color{\qtext}\linespread{1.08}\selectfont #7\par}}};}

% з slides/05-module2.tex
\newcommand{\scopeslot}[2]{%
  \begin{tikzpicture}
    \node[anchor=north west,draw=later,dashed,line width=0.6pt,rounded corners=2pt,inner sep=0.2cm] at (0,0)
      {\topbox{6.55cm}{1.7cm}{{\footnotesize\semibold\color{ink}#1\par}\vskip0.1cm{\tiny\color{mute}#2\par}}};
  \end{tikzpicture}}

% з slides/02-product.tex
\newcommand{\nochip}[1]{\tikz[baseline=(c.base)]\node[draw=line,rounded corners=1pt,inner sep=3pt](c){#1};\enspace}


\title{Lap Coach}
\subtitle{План проєкту на рік}

\begin{document}

% ===== Титул =====
\begin{darkframe}
\begin{frame}[plain,t]
\begin{slidecanvas}
  % шахова клітинка фінішу праворуч угорі
  \checker{(12.4,-3.0)}{9}{7}{0.4}
  \fill[dark,path fading=west] (12.3,-3.1) rectangle ++(1.4,3.2);
  % бордюр унизу
  \kerbstrip{(-0.3,-9)}{34}{0.28}
  % текст
  \node[anchor=north west,align=left,text width=12cm] at (0.9,-1.5) {%
    {\footnotesize\semibold\color{kerb}\caps{Етап 3\enspace·\enspace Планування}}\par\vskip0.3cm
    {\fontsize{46}{46}\selectfont\heavy\color{white}Lap\hspace{0.3em}Coach}\par\vskip0.25cm
    {\large\light\color{white!80!dark}Цифровий тренер для автоспорту}\par\vskip0.7cm
    {\normalsize\semibold\color{white}План проєкту на рік:\ листопад 2026 – жовтень 2027}%
  };
  \node[anchor=south west,align=left,font=\scriptsize,text=later] at (0.9,-8.25)
    {Дарина Демянчук\quad Катерина Наливайко\quad Марина Проніна\quad Ярина Афанасенко\quad Яна Уточкіна};
\end{slidecanvas}
\note{\speaker{Лід · Яна}
Добрий день. Ми команда Lap Coach — цифрового тренера для автоспорту. На етапі 3 ми спланували проєкт на рік: з листопада 2026 до жовтня 2027.

Далі покажемо, що і навіщо ми будуємо, як розклали роботи між людьми, як розставили пріоритети в беклозі і скільки реально встигнемо за модуль 2. Кожен розділ розповідає та, хто його готувала.}
\end{frame}
\end{darkframe}

% ===== Етап у цифрах =====
\begin{frame}{Етап 3 у цифрах}
\vspace{0.15cm}
\begin{columns}[T,onlytextwidth]
  \column{0.31\textwidth}\stat{12}{місяців плану:\\листопад 2026 – жовтень 2027}
  \column{0.31\textwidth}\stat[ink]{25}{робіт у 5 фазах дорожньої карти}
  \column{0.31\textwidth}\stat[ink]{45}{задач у беклозі, B-01…B-45}
\end{columns}
\vspace{0.6cm}
\begin{columns}[T,onlytextwidth]
  \column{0.31\textwidth}\stat[ink]{5}{учасниць, по 6 год на тиждень}
  \column{0.31\textwidth}\stat{48\,год}{на задачі MVP за 2 тижні\\модуля 2, після запасу 20\%}
  \column{0.31\textwidth}\stat[ink]{≤\,0,05\,с}{похибка часу кола — головна\\технічна ціль MVP}
\end{columns}
\vfill
{\scriptsize\color{mute}Артефакти:\enspace Project Overview\enspace·\enspace Roadmap і діаграма Ґанта\enspace·\enspace RACI\enspace·\enspace Беклог\enspace·\enspace Матриця Ейзенхауера\enspace·\enspace Обсяг робіт на модуль 2}
\vspace{0.55cm}
\note{\speaker{Лід · Яна}
Коротко про етап у цифрах. План на 12 місяців. У дорожній карті 25 робіт у п'яти фазах, у беклозі — 45 задач.

Нас п'ятеро, кожна має 6 годин на тиждень. Тому на MVP за два тижні модуля 2 у нас лише 48 годин на задачі — це вже після запасу 20\% на рев'ю, інтеграцію і зустрічі. Це обмеження проходить через увесь план.

Головна технічна ціль MVP — похибка часу кола не більше 0,05 секунди проти офіційного хронометражу.}
\end{frame}

\section{Продукт і цілі}

% ===== Продукт =====
\begin{frame}{Що ми будуємо}
\vspace{0.1cm}
{\small Вебзастосунок, який перетворює GPS-запис з екшн-камери чи телефона на розбір заїзду.\par}

\vspace{0.2cm}
\begin{tikzpicture}[
    pstep/.style={draw=none,fill=soft,minimum height=0.9cm,text width=2.25cm,align=center,font=\scriptsize\semibold,text=ink,inner sep=2pt},
    arr/.style={-{Stealth[length=4pt]},kerb,line width=0.9pt}]
  \node[pstep,fill=ink,text=white] (a) {Відео GoPro\\або трек телефона};
  \node[pstep,right=0.42cm of a] (b) {Поділ на кола};
  \node[pstep,right=0.42cm of b] (c) {Траєкторія на\\межах траси};
  \node[pstep,right=0.42cm of c] (d) {Повороти, де\\втрачено час};
  \node[pstep,right=0.42cm of d,fill=kerb,text=white] (e) {2–3 прості\\поради};
  \foreach \u/\v in {a/b,b/c,c/d,d/e} \draw[arr] (\u) -- (\v);
\end{tikzpicture}

\vspace{0.3cm}
\begin{columns}[T,onlytextwidth]
\column{0.31\textwidth}
  \lbl{Для кого}\par\vskip0.12cm
  {\scriptsize
  \textbf{Прокатний картинг} — пілоти, які їздять регулярно. Головний сегмент MVP.\par\vskip0.08cm
  Далі — аматорські гонки в картингу. Трек-дні на авто — після курсу.\par}
\column{0.31\textwidth}
  \lbl{Проблема}\par\vskip0.12cm
  {\scriptsize
  Новачку бракує об'єктивного зворотного зв'язку.\par\vskip0.08cm
  Тренер — \textbf{€250–350} на день.\\
  Garmin Catalyst 2 — \textbf{\$1,199.99}.\\
  AiM Solo 2 — \textbf{\textasciitilde\$777}.\par}
\column{0.31\textwidth}
  \lbl{Чим відрізняємось}\par\vskip0.12cm
  {\scriptsize
  \textbf{Без нового обладнання:} досить відео GoPro чи треку з телефона, які в пілота вже є.\par\vskip0.08cm
  Поради простою мовою, а не графіки, і відео біля кожної поради.\par}
\end{columns}
\note{\speaker{П1}
Lap Coach — це вебзастосунок-тренер для пілотів-аматорів. Пілот завантажує відео з GoPro або трек з телефона. Застосунок ділить заїзд на кола, накладає траєкторію на межі траси, знаходить повороти, де втрачається час, і дає дві-три прості поради на наступний заїзд.

Головний сегмент MVP — прокатний картинг. Проблема в тому, що новачку бракує зворотного зв'язку: тренер коштує 250–350 євро на день, а телеметрія — від 777 до 1200 доларів.

Наша відмінність — не треба купувати нове обладнання, і поради пишемо простою мовою.}
\end{frame}

% ===== Цілі =====
\begin{frame}{Цілі на курс і на рік}
\vspace{0.05cm}
\centering
\begin{tikzpicture}[x=1.183cm,y=1cm,
    gcard/.style={anchor=north west,inner sep=4pt,fill=soft,font=\tiny,text=body}]
  % дужки «на курс» / «на рік»
  \draw[ink,line width=0.5pt] (0.03,1.42) -- (0.03,1.5) -- (6.97,1.5) -- (6.97,1.42);
  \node[font=\tiny\semibold,text=ink,fill=white,inner sep=3pt] at (3.5,1.5) {На курс · листопад 2026 – травень 2027};
  \draw[ink,line width=0.5pt] (7.03,1.42) -- (7.03,1.5) -- (11.97,1.5) -- (11.97,1.42);
  \node[font=\tiny\semibold,text=ink,fill=white,inner sep=3pt] at (9.5,1.5) {На рік · червень – жовтень 2027};
  % місяці
  \foreach \m [count=\i from 0] in {лис,гру,січ,лют,бер,кві,тра,чер,лип,сер,вер,жов}{
    \node[font=\tiny,text=mute] at (\i+0.5,1.08) {\m};
    \draw[line,line width=0.4pt] (\i,0.88) -- (\i,0.05);
  }
  \draw[line,line width=0.4pt] (12,0.88) -- (12,0.05);
  % фази
  \fill[kerb] (0.04,0.2) rectangle (0.66,0.72);
  \fill[lane] (1.04,0.2) rectangle (6.96,0.72);
  \fill[go]   (8.04,0.2) rectangle (8.96,0.72);
  \fill[ink]  (9.04,0.2) rectangle (10.96,0.72);
  \fill[flag] (11.04,0.2) rectangle (11.96,0.72);
  \node[font=\tiny\semibold,text=white] at (0.35,0.46) {MVP};
  \node[font=\tiny\semibold,text=white] at (4,0.46) {Бета з пілотами};
  \node[font=\tiny\semibold,text=white] at (8.5,0.46) {1.0};
  \node[font=\tiny\semibold,text=white] at (10,0.46) {Більше трас};
  % картки під фазами
  \node[gcard] (m) at (0,-0.15) {\topbox{2.6cm}{3.3cm}{%
     {\scriptsize\semibold\color{kerb}MVP · до \textasciitilde20.11}\par\vskip2pt
     Пілот завантажує відео GoPro з «Жаги швидкості» і бачить кола, повороти з найбільшими втратами і 2–3 поради.\par\vskip4pt
     {\semibold\color{ink}Похибка кола ≤\,0,05\,с}\\ у демо — 0,035\,с}};
  \node[gcard] (b) at (2.45,-0.15) {\topbox{4.25cm}{3.3cm}{%
     {\scriptsize\semibold\color{lane}Бета · грудень – травень}\par\vskip2pt
     Розбори для 5 пілотів, кнопки оцінки порад, логи подій, анкета.\par\vskip4pt
     \renewcommand{\arraystretch}{1.05}%
     \begin{tabular}{@{}>{\semibold\color{ink}}l@{\enspace}>{\raggedright\arraybackslash}p{2.95cm}@{}}
       North Star & ≥\,7 з 10 розборів мають корисну пораду\\
       Активація & ≥\,4 з 5 пілотів доходять до розбору\\
       Утримання & ≥\,40\% — друга сесія за 30 днів\\
       Покращення & ≥\,3 з 5 пілотів, медіана ≥\,0,3\,с
     \end{tabular}}};
  \node[gcard] (r) at (6.42,-0.15) {\topbox{1.7cm}{3.3cm}{%
     {\scriptsize\semibold\color{go}Реліз 1.0}\\{\semibold липень}\par\vskip3pt
     Історія прогресу: найкраще коло по кожній трасі.}};
  \node[gcard] (t) at (8.2,-0.15) {\topbox{1.7cm}{3.3cm}{%
     {\scriptsize\semibold\color{ink}Більше трас}\\{\semibold серпень – вересень}\par\vskip3pt
     Нову відкриту трасу додаємо без ручного обведення меж.}};
  \node[gcard] (f) at (9.98,-0.15) {\topbox{2.15cm}{3.3cm}{%
     {\scriptsize\semibold\color{flag!80!black}Рішення про тарифи}\\{\semibold жовтень}\par\vskip3pt
     Запускаємо freemium, лише якщо в анкеті ≥\,40\% «заплатив би» і утримання ≥\,40\%.}};
  % зв'язки фаза → картка
  \draw[kerb,line width=0.6pt] (0.35,0.2) -- (0.35,-0.15);
  \draw[lane,line width=0.6pt] (4,0.2) -- (4,-0.15);
  \draw[go,line width=0.6pt] (8.5,0.2) -- (8.5,0.02) -- (7.15,0.02) -- (7.15,-0.15);
  \draw[ink,line width=0.6pt] (10,0.2) -- (10,0.02) -- (8.95,0.02) -- (8.95,-0.15);
  \draw[flag,line width=0.6pt] (11.5,0.2) -- (11.5,-0.15);
\end{tikzpicture}
\note{\speaker{П1}
Цілі ми розділили на курс і на рік.

На курс — дві. Перша: MVP до кінця модуля 2, приблизно до 20 листопада. Пілот бачить кола, проблемні повороти і дві-три поради. Міряємо похибкою часу кола: не більше 0,05 секунди, у демо вже вийшло 0,035.

Друга: бета з п'ятьма пілотами з грудня до травня. Головна метрика — North Star: щонайменше 7 з 10 розборів мають хоча б одну корисну пораду.

На рік: у липні Реліз 1.0 з історією прогресу, у серпні–вересні — нові траси без ручної роботи, а в жовтні — рішення про тарифи, тільки якщо бета покаже готовність платити.}
\end{frame}

% ===== Обсяг =====
\begin{frame}{Обсяг: що входить і чого не робимо}
\vspace{0.05cm}
\begin{columns}[T,onlytextwidth]
\column{0.245\textwidth}\phasecol{kerb}{MVP}{%
  Відео GoPro HERO5–HERO11, HERO13\par\vskip2pt
  Автоматичний поділ на кола, без кіл із піт-лейном\par\vskip2pt
  Карта однієї траси, межі вручну за супутником\par\vskip2pt
  Повороти за втратою часу проти найкращого кола\par\vskip2pt
  2–3 поради з цифрами}
\column{0.245\textwidth}\phasecol{lane}{Бета}{%
  Трек із телефона з попередженням про точність\par\vskip2pt
  Фрагмент відео біля поради\par\vskip2pt
  План на наступний заїзд\par\vskip2pt
  Кнопки оцінки порад, логи подій\par\vskip2pt
  Анкета, зокрема про тарифи}
\column{0.245\textwidth}\phasecol{go}{Реліз 1.0}{%
  Історія прогресу: найкраще коло по трасі}
\column{0.245\textwidth}\phasecol{later}{Після курсу}{%
  Автоматичні межі й автопошук поворотів\par\vskip2pt
  Межі за фото з дрона, GPS RaceBox\par\vskip2pt
  Ідеальне коло, «чи виправлено помилку»\par\vskip2pt
  Порівняння з другом, англійська, трек-дні\par\vskip2pt
  Картодроми-партнери, тарифи Pro}
\end{columns}

\vspace{0.3cm}
{\lbl{Свідомо не робимо}}\par\vskip0.12cm
{\tiny\color{body}\nochip{Підказки в реальному часі}\nochip{Криті картодроми}\nochip{Імпорт MyChron і AiM}\nochip{Соцмережа і рейтинги}\nochip{Мобільний застосунок}}
\note{\speaker{П1}
Обсяг розклали за чотирма фазами. У MVP — тільки ядро розбору: відео GoPro, кола, карта однієї траси, повороти з найбільшими втратами і дві-три поради. У бету — трек із телефона, відео біля поради, план на наступний заїзд і все, що потрібно, щоб міряти метрики. Реліз 1.0 — історія прогресу. Решта — після курсу.

І окремо — що ми свідомо не робимо. Підказки в реальному часі — це ніша Garmin Catalyst. На критих картодромах немає GPS. Імпорт MyChron і AiM суперечить нашій ідеї «без нового обладнання». Соцмережа і мобільний застосунок не відповідають на питання «де, чому і що змінити».}
\end{frame}

% ===== Обмеження =====
\begin{frame}{Обмеження}
\vspace{0.15cm}
\begin{columns}[T,onlytextwidth]
\column{0.245\textwidth}\limcard{5 людей × 6 год на тиждень}{ За 2 тижні MVP — 60 год, \textasciitilde48 год на задачі після запасу 20\%}
\column{0.245\textwidth}\limcard{MVP — 4–20 листопада}{Що не влазить у \textasciitilde48 год, переходить у бету}
\column{0.245\textwidth}\limcard{Лише відкриті траси}{На критих картодромах немає GPS}
\column{0.245\textwidth}\limcard{GoPro HERO5–11, HERO13}{У HERO12 немає GPS. Інші камери — трек із телефона, з бети}
\end{columns}
\vspace{0.18cm}
\begin{columns}[T,onlytextwidth]
\column{0.245\textwidth}\limcard{Одна траса в MVP}{«Жага швидкості», Київ: там записано демо. Межі обводимо вручну}
\column{0.245\textwidth}\limcard{Українська мова}{Англійська — після курсу}
\column{0.245\textwidth}\limcard{Бюджет 0}{Безкоштовно для тестувальників. Нове обладнання не купуємо, працюємо з відкритим GPMF}
\column{0.245\textwidth}\limcard{Відео не на сервер}{Дістаємо з файлу лише GPS-телеметрію: кілька МБ на сесію}
\end{columns}
\note{\speaker{П1}
Тепер обмеження. Головне — час. Нас п'ятеро по 6 годин на тиждень, тому за два тижні MVP — 60 годин, а на задачі після запасу — близько 48. Усе, що не влазить, переходить у бету.

Інші обмеження технічні. Працюємо лише на відкритих трасах, бо на критих немає GPS. Підтримуємо GoPro HERO5–HERO11 і HERO13: у HERO12 GPS немає. У MVP одна траса — «Жага швидкості», де записано демо. Інтерфейс український, бюджет нульовий, а відео не вивантажуємо на сервер — тільки телеметрію.}
\end{frame}

% ===== Ризики =====
\begin{frame}{Припущення і ризики}
\vspace{-0.1cm}
\setlength{\aboverulesep}{1.5pt}\setlength{\belowrulesep}{2.5pt}
\renewcommand{\arraystretch}{1.12}
\arrayrulecolor{line}
{\scriptsize
\begin{tabularx}{\textwidth}{@{}>{\raggedright\arraybackslash\leavevmode}p{5.9cm}@{\hspace{0.45cm}}>{\raggedright\arraybackslash\leavevmode}X@{}}
\lbl[mute]{Ризик або припущення} & \lbl[mute]{Що робимо} \\ \midrule
{\semibold\color{ink}Точність GPS GoPro перевірено лише на одній сесії} & Перевіряємо на інших записах у беті. З записів зі слабким сигналом порад не робимо \\ \midrule
{\semibold\color{ink}Точні межі траси потребують фото з дрона; умови знімків Esri не перевірені} & У MVP межі із супутника, невпевнені ділянки — пунктиром. Умови знімків — у перший тиждень модуля 2 \\ \midrule
{\semibold\color{ink}Tenth Finder і Kart Track уже дають поради} & Перевіряємо їх у перший тиждень модуля 2. Наша відмінність — без нового обладнання і з відео біля поради \\ \midrule
{\semibold\color{ink}Виробник може прибрати GPS з камери, як у HERO12} & Публікуємо список камер; для інших — трек із телефона з попередженням про точність \\ \midrule
{\semibold\color{ink}Розмір ринку не досліджено, готовність платити — гіпотеза} & Перші дані дадуть набір пілотів і анкета бети. Тарифи — лише за умов із цілей \\
\end{tabularx}}
\note{\speaker{П1}
П'ять головних ризиків. Найбільший — точність: ми перевірили GPS GoPro лише на одній сесії одного пілота, тому в беті перевіряємо на інших записах. Межі траси точно обвести можна тільки за фото з дрона, а умови супутникових знімків Esri ще не перевірені — це робимо в перший тиждень модуля 2. Тоді ж перевіряємо конкурентів, Tenth Finder і Kart Track.

Ще два ризики — виробник може прибрати GPS з камер, і ми не знаємо розміру ринку. На перший відповідаємо треком із телефона, на другий — анкетою бети.}
\end{frame}

\section{Роботи і люди}

% ===== Дорожня карта + Ґант =====
% Місяці робіт — з gantt.pdf (Roadmap.md), місяць 1 = листопад 2026 (RACI.md).
% TODO: у gantt.pdf роботи MVP (фази 2–4) стоять у місяцях 3–10, а в Project Overview
% MVP — 4–20 листопада, бета — грудень–травень. Коли команда узгодить Ґант
% з Overview, оновити місяці нижче.
\begin{frame}{Дорожня карта: 25 робіт у 5 фазах}
\vspace{-0.1cm}
\centering
\begin{tikzpicture}[x=1.183cm,y=1cm]
  \foreach \m [count=\i] in {лис,гру,січ,лют,бер,кві,тра,чер,лип,сер,вер,жов}{
    \node[font=\scriptsize\semibold,text=ink] at (\i-0.5,0.18) {\i};
    \node[font=\tiny,text=mute] at (\i-0.5,-0.15) {\m};
  }
  \foreach \i in {0,...,12}{\draw[line,line width=0.4pt] (\i,0.38) -- (\i,-4.75);}
  \gphase{1}{2}{Дослідження та підготовка}{1}
  \gphase{3}{5}{Обробка даних (ядро)}{2}
  \gphase{6}{7}{Картографія та аналітика}{3}
  \gphase{8}{10}{Рушій порад та інтерфейс}{4}
  \gphase{11}{12}{Тестування та реліз}{5}
  % Фаза 1
  \gwork{1}{2}{1}{1}{1}{1. Юридичний аудит карт}
  \gwork{1}{2}{1}{2}{2}{2. Проєктування архітектури}
  \gwork{1}{2}{1}{2}{3}{3. UI/UX дизайн}
  \gwork{1}{2}{2}{2}{4}{4. Налаштування CI/CD}
  \gwork{1}{2}{1}{1}{5}{5. Технічна валідація GPMF}
  % Фаза 2
  \gwork{3}{5}{3}{4}{1}{6. Модуль парсингу відео}
  \gwork{3}{5}{3}{4}{2}{7. Алгоритм поділу на кола}
  \gwork{3}{5}{4}{5}{3}{8. Фільтрація невалідних кіл}
  \gwork{3}{5}{4}{5}{4}{9. Корекція зсуву GPS}
  \gwork{3}{5}{3}{5}{5}{10. Серверна інфраструктура}
  % Фаза 3
  \gwork{6}{7}{6}{6}{1}{11. Оцифрування траси MVP}
  \gwork{6}{7}{6}{6}{2}{12. Рендеринг траєкторії}
  \gwork{6}{7}{6}{7}{3}{13. Ідентифікація поворотів}
  \gwork{6}{7}{6}{7}{4}{14. Механізм порівняння кіл}
  \gwork{6}{7}{7}{7}{5}{15. Сортування втрат часу}
  % Фаза 4
  \gwork{8}{10}{8}{9}{1}{16. Рушій текстових порад}
  \gwork{8}{10}{8}{9}{2}{17. Синхронізація з відео}
  \gwork{8}{10}{9}{10}{3}{18. Генератор плану тренувань}
  \gwork{8}{10}{8}{10}{4}{19. Верстка дашборда}
  \gwork{8}{10}{10}{10}{5}{20. Система оцінювання порад}
  % Фаза 5
  \gwork{11}{12}{11}{11}{1}{21. Точнісне тестування}
  \gwork{11}{12}{11}{12}{2}{22. Закрите бета-тестування}
  \gwork{11}{12}{11}{11}{3}{23. Налаштування продуктової аналітики}
  \gwork{11}{12}{12}{12}{4}{24. Валідація бізнес-моделі}
  \gwork{11}{12}{12}{12}{5}{25. Підготовка до масштабування}
\end{tikzpicture}

\vspace{0.05cm}
\raggedright{\tiny\color{mute}Місяць 1 — листопад 2026. Повна діаграма Ґанта — gantt.pdf, опис кожної роботи — Roadmap.md.\par}
\note{\speaker{П2 · Дарина}
Дорожня карта — це 25 робіт у п'яти фазах, які йдуть одна за одною протягом року.

Перші два місяці — дослідження й підготовка: юридичний аудит карт, архітектура, дизайн, CI/CD і перевірка, що бібліотека gopro-telemetry стабільно дістає GPS з відео. Далі три місяці на ядро обробки даних: парсинг відео, поділ на кола, фільтрація і сервер. Потім картографія й аналітика, рушій порад з інтерфейсом, а останні два місяці — тестування точності, закрита бета й аналітика.

Кожна червона смужка — робота на діаграмі Ґанта. Повна діаграма — у файлі gantt.pdf.}
\end{frame}

% ===== RACI =====
% Дані — RACI.md. Рядки — учасниці, стовпці — роботи 1–25 з дорожньої карти.
\begin{frame}{RACI: хто за що відповідає}
\vspace{0.05cm}
\centering
\begin{tikzpicture}[x=1cm,y=1cm]
  % фази і номери робіт
  \foreach \s/\e/\n in {1/5/1,6/10/2,11/15/3,16/20/4,21/25/5}{
    \draw[ink,line width=0.6pt] ({(\s-1)*0.42+0.03},0.78) -- ({\e*0.42-0.03},0.78);
    \node[font=\tiny\semibold,text=ink] at ({(\s+\e-1)*0.21},0.92) {Ф\n};
  }
  \foreach \c in {1,...,25}{\node[font=\tiny\ttfamily,text=mute] at ({(\c-1)*0.42+0.21},0.5) {\c};}
  \node[font=\tiny\semibold,text=ink] at (10.95,0.5) {R};
  \node[font=\tiny\semibold,text=kerb] at (11.5,0.5) {A};
  \racirow{1}{Демянчук Д.}{бек/тест}{I,AR,I,R,AR,C,C,R,C,AR,I,I,AR,C,I,R,C,C,R,AR,C,C,C,I,R}
  \racirow{2}{Наливайко К.}{аналітик/бек}{C,C,C,I,I,I,AR,C,AR,I,C,C,R,AR,C,AR,I,C,I,I,R,C,AR,C,R}
  \racirow{3}{Проніна М.}{фронт}{I,C,AR,I,C,AR,I,-,-,C,R,AR,I,-,-,I,AR,I,AR,C,I,C,R,C,C}
  \racirow{4}{Афанасенко Я.}{бек/тест}{I,C,I,AR,C,I,R,AR,C,C,I,I,C,R,AR,I,R,AR,I,R,AR,R,I,I,C}
  \racirow{5}{Уточкіна Я.}{лід}{AR,I,R,I,I,I,I,I,I,I,AR,I,I,-,-,R,I,R,C,C,I,AR,I,AR,A}
  % підсумки
  \foreach \r/\R in {1/10,2/8,3/7,4/10,5/7}{
    \node[font=\scriptsize\heavy,text=ink] at (10.95,{-(\r-1)*0.62}) {\R};
    \node[font=\scriptsize\heavy,text=kerb] at (11.5,{-(\r-1)*0.62}) {5};
  }
  \draw[line] (10.65,0.35) -- (10.65,-2.8);
\end{tikzpicture}

\vspace{0.25cm}
{\tiny\color{body}
\tikz[baseline=-0.6ex]\fill[ink] (0,-0.1) rectangle (0.2,0.12);\;\textbf{R} — виконує\qquad
\tikz[baseline=-0.6ex]\fill[kerb] (0,-0.1) rectangle (0.2,0.12);\;\textbf{A/R} — виконує і відповідає\qquad
\tikz[baseline=-0.6ex]\draw[kerb,line width=0.8pt] (0,-0.1) rectangle (0.2,0.12);\;\textbf{A} — відповідає, рівно одна на роботу\qquad
\tikz[baseline=-0.6ex]\fill[line] (0,-0.1) rectangle (0.2,0.12);\;\textbf{C} — радяться\qquad
{\color{later}\textbf{I}} — інформують}
\note{\speaker{П4 · Ярина}
RACI-матриця: рядки — це ми п'ятеро, стовпці — 25 робіт з дорожньої карти, з тими самими номерами.

Правила такі: у кожній роботі рівно одна людина відповідає за результат — це A, і хоча б одна виконує — це R. Червоні клітинки — коли одна людина і виконує, і відповідає.

Праворуч — підсумки. Кожна з нас відповідає за п'ять робіт. Виконавицею найчастіше стоять Дарина і Ярина — по 10 робіт, Катерина — 8, Марина і Яна — по 7.}
\end{frame}

% ===== Навантаження =====
% Дані — RACI.md, розділ «Навантаження по місяцях».
\begin{frame}{Навантаження по місяцях}
\vspace{0.1cm}
\begin{columns}[T,onlytextwidth]
\column{0.68\textwidth}
\begin{tikzpicture}
  \foreach \m [count=\i] in {лис,гру,січ,лют,бер,кві,тра,чер,лип,сер,вер,жов}{
    \node[font=\tiny\ttfamily,text=mute] at ({(\i-0.5)*0.62},0.62) {\i};
    \node[font=\tiny,text=mute] at ({(\i-0.5)*0.62},0.4) {\m};}
  \loadrow{1}{Демянчук Д.}{2,2,1,2,2,1,1,2,2,2,0,1}
  \loadrow{2}{Наливайко К.}{0,0,1,2,1,2,2,1,1,0,2,1}
  \loadrow{3}{Проніна М.}{1,1,1,1,0,2,0,2,2,1,1,0}
  \loadrow{4}{Афанасенко Я.}{0,1,1,2,1,1,2,1,2,2,2,1}
  \loadrow{5}{Уточкіна Я.}{2,1,0,0,0,1,0,1,2,1,1,2}
\end{tikzpicture}

\vspace{0.2cm}
{\tiny\color{mute}Кількість одночасних робіт, де людина R. Місяць 1 — листопад 2026.\par}
\column{0.29\textwidth}
  \lbl{Перевірка матриці}\par\vskip0.2cm
  {\scriptsize
  \ok У кожній роботі рівно одна A\par\vskip0.12cm
  \ok У кожній роботі хоча б одна R\par\vskip0.12cm
  \ok R пораховано на кожну\par\vskip0.12cm
  \ok Ніхто не має більше 2 R одночасно}
\end{columns}
\note{\speaker{П4 · Ярина}
Щоб ніхто не перевантажився, ми розклали RACI по місяцях за діаграмою Ґанта. У клітинці — скільки робіт людина виконує одночасно в цьому місяці.

Максимум — дві роботи одночасно, більше ні в кого немає. Ми також перевірили, що в кожній роботі рівно одна A і хоча б одна R, і порахували R на кожну учасницю.}
\end{frame}

\section{Беклог і пріоритети}

% ===== Беклог =====
% Дані — Backlog.pdf, розділи 1–2. Пріоритети пораховано за рядками беклогу.
\begin{frame}{Беклог: 45 задач за фазами}
\vspace{0.15cm}
\centering
\begin{tikzpicture}
  \node[anchor=east,font=\tiny\semibold,text=mute] at (14.2,0.62) {\caps{Оцінка, год}};
  \node[anchor=west,font=\tiny\semibold,text=mute] at (2.9,0.62) {\caps{Задач за пріоритетом MoSCoW}};
  \brow{1}{MVP}{4–20 листопада 2026}{19}{5}{0}{68}
  \brow{2}{Бета}{грудень 2026 – травень 2027}{11}{1}{0}{51}
  \brow{3}{Реліз 1.0}{липень 2027}{2}{0}{0}{8}
  \brow{4}{Після курсу}{з червня 2027}{0}{2}{5}{—}
  \draw[line] (0,-2.72) -- (14.2,-2.72);
  \node[anchor=west,font=\footnotesize\semibold,text=ink] at (0,-3.0) {Разом 45};
  \node[anchor=west,font=\tiny,text=body] at (2.9,-3.0) {%
    \tikz[baseline=-0.6ex]\fill[kerb] (0,-0.1) rectangle (0.22,0.12);\;Must 32\qquad
    \tikz[baseline=-0.6ex]\fill[kerb!35] (0,-0.1) rectangle (0.22,0.12);\;Should 8\qquad
    \tikz[baseline=-0.6ex]\fill[line] (0,-0.1) rectangle (0.22,0.12);\;Won't 5};
  \node[anchor=east,font=\tiny,text=mute,align=right] at (14.2,-3.0) {MVP: Must 56 + Should 12\\Бета: Must 47 + Should 4};
\end{tikzpicture}

\vspace{0.15cm}
\raggedright
{\tiny\color{mute}\caps{Епіки}\enspace
\color{body}F1 Завантаження заїзду\enspace·\enspace F2 Поділ на кола\enspace·\enspace F3 Карта траси і повороти\enspace·\enspace F4 Де втрачається час\enspace·\enspace F5 Поради простою мовою\enspace·\enspace F6 План на наступний заїзд\enspace·\enspace F7 Історія прогресу\enspace·\enspace Технічна основа\enspace·\enspace Перевірки\enspace·\enspace Бета і метрики\par}
\note{\speaker{П3}
Беклог — це 45 задач з номерами від B-01 до B-45. Кожна має фазу, пріоритет за MoSCoW, оцінку в годинах і залежності. Епіки F1–F7 — це функції з Product Map, плюс технічна основа, перевірки, бета і метрики.

Найбільше задач у MVP — 24, з них 19 обов'язкових, Must. Разом вони тягнуть на 68 годин. У беті 12 задач на 51 годину, у Реліз 1.0 — дві задачі на 8 годин. Сім задач після курсу ми поки не оцінюємо.}
\end{frame}

% ===== MVP: ресурс і обсяг =====
% Дані — Backlog.pdf, розділ 3 «Обсяг MVP і ресурс».
\begin{frame}{MVP: обсяг проти ресурсу}
\vspace{0.1cm}
\begin{columns}[T,onlytextwidth]
\column{0.6\textwidth}
\begin{tikzpicture}
  \node[anchor=west,font=\scriptsize\semibold,text=ink] at (0,0) {Ресурс};
  \fill[ink] (2.2,-0.22) rectangle ({2.2+48*\hb},0.22);
  \fill[line] ({2.2+48*\hb},-0.22) rectangle ({2.2+60*\hb},0.22);
  \node[font=\tiny\semibold,text=white] at ({2.2+24*\hb},0) {48 год на задачі};
  \node[font=\tiny\semibold,text=mute] at ({2.2+54*\hb},0) {+12};
  \node[anchor=west,font=\tiny,text=mute] at ({2.2+60*\hb+0.05},0) {запас};
  %
  \node[anchor=west,font=\scriptsize\semibold,text=ink] at (0,-0.85) {Беклог MVP};
  \fill[kerb] (2.2,-1.07) rectangle ({2.2+56*\hb},-0.63);
  \fill[kerb!35] ({2.2+56*\hb},-1.07) rectangle ({2.2+68*\hb},-0.63);
  \node[font=\tiny\semibold,text=white] at ({2.2+28*\hb},-0.85) {Must 56};
  \node[font=\tiny\semibold,text=ink] at ({2.2+62*\hb},-0.85) {12};
  %
  \node[anchor=west,font=\scriptsize\semibold,text=ink] at (0,-1.7) {Після спрощень};
  \fill[kerb] (2.2,-1.92) rectangle ({2.2+48*\hb},-1.48);
  \node[font=\tiny\semibold,text=white] at ({2.2+24*\hb},-1.7) {Must 48};
  %
  \draw[kerb,dashed,line width=0.7pt] ({2.2+48*\hb},0.45) -- ({2.2+48*\hb},-2.15);
  \node[anchor=north,font=\tiny\semibold,text=kerb] at ({2.2+48*\hb},-2.15) {48};
  \draw[mute,line width=0.4pt] (2.2,-2.15) -- (2.2,0.45);
\end{tikzpicture}

\vspace{0.25cm}
{\scriptsize\color{body}
\textbf{Ресурс:} 2 тижні × 6 год × 5 людей = 60 год, мінус запас 20\% на рев'ю, інтеграцію й зустрічі.\par\vskip0.1cm
\textbf{Перевищення:} Must — 56 год, це на 8 год більше. Should (12 год) переходять у бету першими.\par}
\column{0.36\textwidth}
  \lbl{Як вкластися: −8 год}\par\vskip0.15cm
  \renewcommand{\arraystretch}{1.35}
  {\tiny
  \begin{tabularx}{\linewidth}{@{}>{\ttfamily}lXr@{}}
  B-05 & Лише виїзне коло, без піт-лейну & 3→1 \\
  B-08 & Карта — у бету, у MVP список поворотів & 3→0 \\
  B-13, B-17 & Один екран розбору замість двох & 5→3 \\
  B-09 & Без швидкості на виході & 5→4 \\
  \arrayrulecolor{line}\midrule
  \multicolumn{2}{@{}l}{\semibold Разом Must} & \semibold 56→48 \\
  \end{tabularx}}
  \vskip0.2cm
  {\tiny\color{mute}Ціль MVP не змінюється: кола, повороти з найбільшими втратами і 2–3 поради.\par}
\end{columns}
\note{\speaker{П3}
Тепер найважливіше — чи влазить MVP у наш час. Ресурс: два тижні, по 6 годин, п'ятеро людей — 60 годин. Мінус 20\% на рев'ю, інтеграцію і зустрічі, тож на задачі лишається 48.

А обов'язкові задачі MVP — це 56 годин, на вісім більше. Тому, по-перше, усі задачі Should переходять у бету першими. По-друге, ми спрощуємо чотири задачі: відкидаємо лише виїзне коло, карту переносимо в бету, робимо один екран розбору замість двох і не рахуємо швидкість на виході. Виходить рівно 48 годин, а ціль MVP не змінюється.}
\end{frame}

% ===== Критичний шлях =====
% Дані — Backlog.pdf, розділ 3 «Порядок робіт».
\begin{frame}{Критичний шлях MVP: 28 годин}
\vspace{0.3cm}
\centering
\begin{tikzpicture}
  \def\hw{0.455}% см на годину
  \def\x{0}
  \foreach \idd/\h/\name [count=\i] in {
      B-20/3/Перевірка gopro-telemetry,
      B-04/4/Час кожного кола,
      B-10/4/Відрізки між поворотами,
      B-09/5/Гальмування і швидкість у повороті,
      B-14/5/Правила порад,
      B-15/3/До 3 порад з цифрами,
      B-17/2/Екран розбору,
      B-22/2/Прогін демо}{%
    \pgfmathsetmacro{\w}{\h*\hw}
    \ifodd\i\def\fc{ink}\else\def\fc{kerb}\fi
    \fill[\fc] (\x,0) rectangle ++(\w-0.08,0.75);
    \node[font=\tiny\ttfamily\semibold,text=white,anchor=north west] at (\x+0.05,0.7) {\idd};
    \node[font=\tiny\semibold,text=white,anchor=south west] at (\x+0.05,0.03) {\h\,год};
    \node[font=\tiny,text=body,anchor=north west,text width=\w cm-0.12cm,align=left,inner sep=0pt]
      at (\x,-0.12) {\name};
    \pgfmathsetmacro{\nx}{\x+\w}\global\let\x\nx
  }
  \draw[decorate,decoration={brace,amplitude=4pt,mirror},mute] (0,-1.05) -- (12.74,-1.05);
  \node[font=\scriptsize\semibold,text=ink] at (6.37,-1.4) {28 год з 48 — решту задач робимо паралельно};
\end{tikzpicture}

\vspace{0.55cm}
\raggedright
\lbl{Старт без залежностей}\par\vskip0.15cm
\begin{columns}[T,onlytextwidth]
\column{0.24\textwidth}{\scriptsize\id{B-18}\par\color{body}\tiny Каркас вебзастосунку, збирання, CI, деплой\par}
\column{0.24\textwidth}{\scriptsize\id{B-19}\par\color{body}\tiny Формат даних: JSON сесії і конфігурація траси\par}
\column{0.24\textwidth}{\scriptsize\id{B-20}\par\color{body}\tiny Перевірка gopro-telemetry на відео демо, тестові дані\par}
\column{0.24\textwidth}{\scriptsize\id{B-23}\par\color{body}\tiny Умови використання знімків Esri\par}
\end{columns}
\note{\speaker{П3}
Критичний шлях — це ланцюжок задач, кожна з яких чекає на попередню. У нас він від перевірки бібліотеки gopro-telemetry до прогону демо: вісім задач, 28 годин. Ширина кожного блоку — її оцінка в годинах.

Почати можна з чотирьох задач без залежностей: каркас застосунку, формат даних, перевірка gopro-telemetry і умови знімків Esri. Коли є формат даних і тестові дані з демо, поділ на кола, повороти й поради можна робити паралельно із завантаженням відео.}
\end{frame}

% ===== Ейзенхауер =====
% Дані — Eisenhower.pdf.
\begin{frame}{Матриця Ейзенхауера}
\vspace{-0.2cm}
\begin{columns}[T,onlytextwidth]
\column{0.67\textwidth}
\begin{tikzpicture}
  \node[font=\tiny\semibold,text=mute] at (2.55,0.22) {\caps{Терміново}};
  \node[font=\tiny\semibold,text=mute] at (7.0,0.22) {\caps{Не терміново}};
  \node[font=\tiny\semibold,text=mute,rotate=90] at (-0.22,-1.26) {\caps{Важливо}};
  \node[font=\tiny\semibold,text=mute,rotate=90] at (-0.22,-3.7) {\caps{Не важливо}};
  \def\qtext{white}\def\qsub{white!75!kerb}
  \qcell{0.35}{0}{kerb}{Q1 · Робимо зараз}{17}{ядро MVP, 47 год після спрощень}{\id{B-01} відео GoPro, GPS у браузері\\ \id{B-04} час кожного кола\\ \id{B-11} втрата часу в повороті\\ \id{B-15} до 3 порад з цифрами\\ \id{B-21} автотест точності}
  \def\qtext{white}\def\qsub{white!70!ink}
  \qcell{4.83}{0}{ink}{Q2 · Плануємо}{18}{бета і реліз 1.0}{\id{B-08} карта траси\\ \id{B-25} GPX з телефона\\ \id{B-29} сервер і база\\ \id{B-32} оцінка порад (North Star)\\ \id{B-37} список сесій}
  \def\qtext{ink}\def\qsub{mute}
  \qcell{0.35}{-2.48}{soft}{Q3 · Швидко й мінімально}{1}{потрібна до демо, 1 год}{\id{B-23} умови знімків Esri — у перший тиждень, бо блокує межі траси (B-06)}
  \qcell{4.83}{-2.48}{soft}{Q4 · Відкладаємо}{9}{після курсу або якщо лишиться час}{\id{B-24} Tenth Finder і Kart Track\\ \id{B-35} анкета бети\\ \id{B-39} більше трас\\ \id{B-41} ідеальне коло}
\end{tikzpicture}
\column{0.3\textwidth}
  \vspace{0.3cm}
  \lbl{Критерії}\par\vskip0.12cm
  {\tiny\color{body}
  \textbf{Важливо} — впливає на North Star або без неї не працює розбір.\par\vskip0.08cm
  \textbf{Терміново} — потрібна для MVP (до \textasciitilde20.11) або блокує таку задачу.\par}
  \vskip0.35cm
  \lbl{Звірка з MVP}\par\vskip0.12cm
  {\tiny\color{body}
  \textbf{Q1 + Q3 = MVP:} 18 задач, 48 год — рівно ресурс команди.\par\vskip0.08cm
  Усі задачі розкладено:\\ 17 + 18 + 1 + 9 = 45.\par}
\end{columns}
\note{\speaker{П4 · Ярина}
Матриця Ейзенхауера розкладає всі 45 задач беклогу на чотири квадранти. Важливо — якщо задача впливає на North Star або без неї не працює розбір. Терміново — якщо вона потрібна для MVP або блокує таку задачу.

У першому квадранті 17 задач — це і є ядро MVP. У другому 18: їх плануємо на бету і реліз. У третьому одна — перевірити умови знімків Esri, бо від цього залежать межі траси. У четвертому дев'ять задач, які відкладаємо.

Перший і третій квадранти разом — 18 задач на 48 годин, рівно наш ресурс на MVP.}
\end{frame}

\section{Модуль 2}

% Дані розділу — Semester_Scope.md (обсяг робіт на модуль 2, автор — П5).

% ===== Ціль MVP і що переходить у бету =====
% Semester_Scope.md, розділи 1, 3 і 4.
\begin{frame}{Обсяг робіт на модуль 2}
\vspace{0.05cm}
\lbl{Ціль MVP}\par\vskip0.15cm
\begin{tikzpicture}
  \node[anchor=north west,inner sep=0pt,text width=13.9cm,align=left] (g) at (0.3,0)
    {{\footnotesize\semibold\color{ink}\linespread{1.2}\selectfont Пілот завантажує відео GoPro з «Жаги швидкості» і в~браузері бачить свої кола, повороти, де втратив найбільше часу, і~2–3 поради з~цифрами; час кола відрізняється від офіційного хронометражу в~середньому не більше ніж на 0,05\,с.\par}};
  \fill[kerb] (0,0) rectangle ($(g.south west)+(-0.24,0)$);
\end{tikzpicture}

\vspace{0.2cm}
\begin{columns}[T,onlytextwidth]
\column{0.17\textwidth}\stat{18}{задач у MVP на \textbf{48 год} — рівно ресурс}
\column{0.17\textwidth}\stat[ink]{−8}{год: 4 спрощення Must — \id{B-05}, \id{B-08}, \id{B-09}, \id{B-13}}
\column{0.17\textwidth}\stat[lane]{15}{год: 6 задач переходять у бету}
\column{0.43\textwidth}
  \lbl[lane]{Переходить у бету}\par\vskip0.1cm
  \renewcommand{\arraystretch}{1.12}
  {\tiny
  \begin{tabularx}{\linewidth}{@{}>{\ttfamily}lXr@{}}
  B-08 & Карта траси з траєкторією (Must) & 3 \\
  B-03 & Повідомлення для камер без GPS & 2 \\
  B-06 & Межі траси за супутником & 4 \\
  B-12 & Відкидання кіл із трафіком & 2 \\
  B-16 & Мітка сили зв'язку біля поради & 2 \\
  B-24 & Перевірка Tenth Finder і Kart Track & 2 \\
  \end{tabularx}}
  \vskip0.05cm
  {\tiny\color{mute}Окрім \id{B-08}, усі — Should: вони переходять першими.\par}
\end{columns}
\note{\speaker{П5}
Ціль MVP одним реченням: пілот завантажує відео GoPro з «Жаги швидкості» і в браузері бачить свої кола, повороти, де втратив найбільше часу, і 2–3 поради з цифрами. Час кола має відрізнятися від офіційного хронометражу в середньому не більше ніж на 0,05 секунди.

У MVP — 18 задач на 48 годин, рівно наш ресурс. Щоб вкластися, ми спростили чотири обов'язкові задачі й зекономили 8 годин. Шість задач на 15 годин переходять у бету: це п'ять задач Should і карта траси.}
\end{frame}

% ===== Розподіл по людях =====
% Semester_Scope.md, розділ 5. Червоним — задачі критичного шляху (розділ 5, «Порядок»).
\begin{frame}{Розподіл по людях: по 12 годин}
\vspace{0.2cm}
\centering
\begin{tikzpicture}
  \node[anchor=west,font=\tiny\semibold,text=mute] at (3.3,0.62) {\caps{Задачі MVP, ширина — оцінка в годинах}};
  \node[anchor=east,font=\tiny\semibold,text=mute] at (14.2,0.62) {\caps{Разом, год}};
  \pplrow{1}{У1}{Платформа і завантаження}{B-18/3/ink,B-01/4/ink,B-02/2/ink}{3}
  \pplrow{2}{У2 · лід}{Дані, точність, демо}{B-19/2/ink,B-20/3/kerb,B-21/2/ink,B-23/1/ink,B-22/2/kerb}{2}
  \pplrow{3}{У3}{Кола і відрізки}{B-04/4/kerb,B-05/1/ink,B-10/4/kerb}{3}
  \pplrow{4}{У4}{Повороти і втрати}{B-07/2/ink,B-09/4/kerb,B-11/3/ink,B-13/1/ink}{2}
  \pplrow{5}{У5}{Поради й екран розбору}{B-14/5/kerb,B-15/3/kerb,B-17/2/kerb}{2}
  \draw[line] (0,-3.88) -- (14.2,-3.88);
  \node[anchor=west,font=\footnotesize\semibold,text=ink] at (0,-4.18) {18 задач};
  \node[anchor=west,font=\tiny,text=body] at (3.3,-4.18) {%
    \tikz[baseline=-0.6ex]\fill[kerb] (0,-0.1) rectangle (0.22,0.12);\;критичний шлях, 27 год\qquad
    \tikz[baseline=-0.6ex]\fill[ink] (0,-0.1) rectangle (0.22,0.12);\;інші задачі\qquad
    \tikz[baseline=-0.6ex]\fill[soft] (0,-0.1) rectangle (0.22,0.12);\;запас};
  \node[anchor=east,font=\tiny,text=mute] at (14.2,-4.18) {48 на задачі + 12 запас = 60};
\end{tikzpicture}
\note{\speaker{П5}
У команді п'ять потоків, кожен — приблизно 12 годин за два тижні: 9–10 годин на задачі і 2–3 години запасу. Ширина блоку — оцінка задачі в годинах.

Перший учасник робить платформу і завантаження відео. Лід — дані, точність і демо: формат даних, тестові дані і фінальний прогін зводять роботу всіх. Третій — поділ на кола і відрізки, четвертий — повороти і втрати часу, п'ятий — поради й екран розбору.

Червоним позначено критичний шлях. У MVP він 27 годин, бо метрики повороту ми спростили на годину.}
\end{frame}

% ===== Порядок робіт =====
% Semester_Scope.md, розділ 5 «Порядок». Відкриті питання — Project-Overview.md і звірка артефактів.
\begin{frame}{Порядок на два тижні}
\vspace{0cm}
\centering
\begin{tikzpicture}
  \def\dw{0.8875}% см на день, 4–20 листопада
  % тижні
  \fill[ink] (0,0) rectangle ({7*\dw-0.04},0.42);
  \fill[kerb] ({8*\dw},0) rectangle ({16*\dw},0.42);
  \node[anchor=west,font=\tiny\semibold,text=white] at (0.1,0.21) {ТИЖДЕНЬ 1\enspace·\enspace 4–11 листопада};
  \node[anchor=west,font=\tiny\semibold,text=white] at ({8*\dw+0.1},0.21) {ТИЖДЕНЬ 2\enspace·\enspace 12–20 листопада};
  % перші 2 дні
  \draw[mute,line width=0.5pt] (0,0.55) -- (0,0.75) -- ({2*\dw},0.75) -- ({2*\dw},0.55);
  \node[anchor=south west,font=\tiny,text=body,inner sep=1pt] at (0,0.78) {перші 2 дні: \id{B-18}, \id{B-19}, \id{B-20}};
  % дедлайн прогону
  \draw[kerb,dashed,line width=0.7pt] ({14*\dw},-0.1) -- ({14*\dw},0.95);
  \node[anchor=south east,font=\tiny,text=mute,inner sep=1pt] at ({16*\dw},0.78) {{\semibold\color{kerb}\id{B-22} до 18.11}\enspace·\enspace 2 дні на виправлення};
  % що в тижнях
  \node[anchor=north west,font=\tiny,text=body,text width=6.0cm,align=left,inner sep=0pt] at (0,-0.15)
    {\id{B-23} — у перший тиждень, як вимагає Project Overview.\\ Паралельно: \id{B-01}, \id{B-02}; \id{B-04}, \id{B-05}; \id{B-07}; \id{B-21} після \id{B-04}.};
  \node[anchor=north west,font=\tiny,text=body,text width=6.9cm,align=left,inner sep=0pt] at ({8*\dw},-0.15)
    {\id{B-10} → \id{B-09}, \id{B-11}, \id{B-13} → \id{B-14}, \id{B-15}, \id{B-17} → \id{B-22}};
\end{tikzpicture}

\vspace{0.2cm}
\raggedright
\begin{columns}[T,onlytextwidth]
\column{0.47\textwidth}
  \lbl{Критичний шлях: 27 год}\par\vskip0.12cm
  {\tiny\color{body}При 6 год на тиждень на людину він не влазить у 2 тижні, якщо чекати готового коду попередника.\par\vskip0.1cm
  Тому кожен модуль пишемо і тестуємо на тестових даних демо (\id{B-20}) у форматі \id{B-19}. У4 і У5 стартують у першому тижні з вручну підготовленого JSON сесії. Зводимо все в \id{B-22}.\par}
\column{0.47\textwidth}
  \lbl[ink]{Відкриті питання}\par\vskip0.12cm
  {\tiny\color{body}
  \begin{itemize}\setlength{\itemsep}{0.06cm}
    \item \textbf{Технології не обрані.} Поки припускаємо TypeScript у браузері, gopro-telemetry і MVP без сервера.
    \item \textbf{Tenth Finder і Kart Track.} Overview перевіряє їх у перший тиждень, а в обсязі MVP \id{B-24} — у беті.
    \item \textbf{Тривалість.} 4–20 листопада — це 2,4 тижня, а не 2. Якщо рахувати весь період, на задачі буде 58 год.
  \end{itemize}}
\end{columns}
\note{\speaker{Лід · Яна}
Перші два дні — каркас, формат даних і перевірка gopro-telemetry: від них залежать усі. Прогін демо — до 18 листопада, щоб лишилося два дні на виправлення.

Критичний шлях — 27 годин: якщо чекати код попередника, він не влазить. Тому кожен тестує свій модуль на даних демо у спільному форматі, а зводимо все в кінці.

Відкриті питання — на слайді: технології, перевірка конкурентів і тривалість модуля.}
\end{frame}

% ===== Definition of Done =====
% Semester_Scope.md, розділ 6.
\begin{frame}{Definition of Done}
\vspace{0.2cm}
\begin{columns}[T,onlytextwidth]
\column{0.47\textwidth}
  \lbl[ink]{Задача готова, коли}\par\vskip0.25cm
  {\scriptsize\color{body}
  \dod{Код у головній гілці, пройшов рев'ю одного учасника команди.}
  \dod{Є автотести на тестових даних демо (\id{B-20}); CI зелений.}
  \dod{Вхід і вихід модуля відповідають формату \id{B-19}.}
  \dod{Тексти для пілота — українською, за tone of voice з Product Map.}
  \dod{Відео не покидає комп'ютер пілота: на мережу йде лише код застосунку.}
  \dod{Задеплоєно на хостинг (\id{B-18}), якщо задача видна користувачу.}}
\column{0.47\textwidth}
  \lbl{MVP готовий, коли}\par\vskip0.25cm
  {\scriptsize\color{body}
  \dod{Усі 18 задач MVP відповідають DoD задачі.}
  \dod{Автотест \id{B-21} проходить: похибка часу кола проти AMB у середньому ≤\,0,05\,с на демо-сесії.}
  \dod{Сценарій демо проходить від початку до кінця на задеплоєній версії в Chrome на ноутбуці без помилок у консолі.}
  \dod{Задачі, що переїжджають у бету, внесені в беклог бети.}}
\end{columns}
\note{\speaker{П5}
Definition of Done — коли ми вважаємо роботу завершеною. Задача готова, коли код у головній гілці після рев'ю, є автотести на даних демо і CI зелений, модуль працює у спільному форматі даних, тексти українською, відео не покидає комп'ютер пілота, а все видиме користувачу задеплоєно.

MVP готовий, коли всі 18 задач виконані за цим правилом, автотест точності проходить, сценарій демо проходить на задеплоєній версії без помилок, а все, що переїжджає, внесено в беклог бети.}
\end{frame}

% ===== Сценарій демо =====
% Semester_Scope.md, розділ 7.
\begin{frame}{Сценарій демо: \textasciitilde5 хвилин}
\vspace{-0.05cm}
{\scriptsize\color{mute}Сесія «Жага швидкості», 1 жовтня 2026\enspace·\enspace GoPro HERO7, \id{GH011616} і \id{GH021616}\enspace·\enspace 14 кіл\par}
\vspace{0.2cm}
\centering
\begin{tikzpicture}
  \demostep{0}{0}{1}{Відкриваємо застосунок за посиланням}{Екран завантаження.}
  \demostep{3.6}{0}{2}{Обираємо обидва файли, показуємо вкладку мережі}{Одна сесія з двох частин; відео не вивантажується.}
  \demostep{7.2}{0}{3}{Таблиця кіл}{14 залікових кіл, виїзне відкинуто; найкраще — коло 11, 41,474\,с.}
  \demostep{10.8}{0}{4}{Результат автотесту \id{B-21}}{Похибка проти AMB у середньому ≤\,0,05\,с (у демо — 0,035\,с).}
  \demostep{0}{-2.12}{5}{Список поворотів T1–T7}{Відсортовано за втраченим часом, найбільші втрати першими.}
  \demostep{3.6}{-2.12}{6}{Картки порад}{2–3 поради з цифрою «зараз» і ціллю.}
  \demostep{7.2}{-2.12}{7}{Клікаємо пораду}{Перехід до цього повороту у списку.}
  \demostep{10.8}{-2.12}{8}{Що далі в беті}{Карта траси, відео біля поради, трек із телефона, оцінка порад.}
\end{tikzpicture}

\vspace{0.15cm}
\raggedright
{\tiny\color{body}\textbf{\color{kerb}Запасний план:} якщо завантаження не працює на демо-ноутбуці — відкриваємо збережений JSON сесії (\id{B-20}) і показуємо кроки 3–7.\par}
\note{\speaker{П5}
Демо MVP — приблизно п'ять хвилин. Беремо реальну сесію з «Жаги швидкості» — два файли GoPro, 14 кіл.

Обираємо обидва файли і показуємо, що відео не вивантажується. Далі — таблиця кіл, автотест точності, повороти за втраченим часом і картки порад; клік на пораду веде до повороту. Наприкінці — що буде в беті.

Якщо не спрацює завантаження, відкриваємо збережений JSON сесії.}
\end{frame}

% ===== Фінал =====
\begin{darkframe}
\begin{frame}[plain,t]
\begin{slidecanvas}
  \kerbstrip{(-0.3,-9)}{34}{0.28}
  \checker{(12.4,-3.0)}{9}{7}{0.4}
  \fill[dark,path fading=west] (12.3,-3.1) rectangle ++(1.4,3.2);
  \node[anchor=north west,align=left,text width=9cm] at (0.9,-1.6) {%
    {\fontsize{40}{40}\selectfont\heavy\color{white}Дякуємо}\par\vskip0.3cm
    {\large\light\color{white!80!dark}Питання?}};
  \node[anchor=north west,align=left,font=\scriptsize,text=later,text width=11cm] at (0.9,-4.3) {%
    {\semibold\color{kerb}\caps{Артефакти етапу 3}}\par\vskip0.15cm
    {\ttfamily docs/stages/stage-3/}\par\vskip0.08cm
    {\color{white}Project-Overview.md}\enspace опис, цілі, обсяг, обмеження, ризики\\
    {\color{white}Roadmap.md, gantt.pdf}\enspace перелік робіт і діаграма Ґанта\\
    {\color{white}RACI.md}\enspace розподіл відповідальності і навантаження\\
    {\color{white}Backlog.pdf}\enspace беклог продукту\\
    {\color{white}Eisenhower.pdf}\enspace матриця Ейзенхауера};
\end{slidecanvas}
\note{\speaker{Лід · Яна}
Це наш план на рік. Усі артефакти лежать у репозиторії, у папці docs/stages/stage-3. Дякуємо за увагу — будемо раді питанням.}
\end{frame}
\end{darkframe}


\end{document}
